\documentclass{article}
\usepackage[utf8]{inputenc}
\usepackage{amsmath}
\usepackage{amssymb}
\usepackage{graphicx}
\usepackage{tabularx}
\usepackage{caption}
\usepackage{booktabs}
\usepackage[english]{babel}
\usepackage{setspace}
\usepackage{geometry}
\usepackage{float}
\usepackage[justification=centering]{caption}
\usepackage{ragged2e}

% APA style formatting
\geometry{a4paper, margin=1in}
\setstretch{1.5}
\captionsetup[table]{labelfont=bf, justification=raggedright, singlelinecheck=false}
\captionsetup[figure]{font=small, labelfont=bf, justification=raggedright, singlelinecheck=false}

\title{Cloud Native Streaming Platform Netflix-inspired}
\author{ Santiago Antonio Gualteros Bustos \\ Nicolas Castro Rivera \\ Andrés \\ Erick Santiago Buitrago Peña   }
\date{}

\begin{document}

\maketitle

\section{Executive Summary}
Netflix is a streaming platform that provides this service for subscriptions. Whose vision is entertaining customers through series and movies.
This content can be accessed around the world by millions of customers those who want to watch their favorites shows at any time.
For that reason is an application demands great aviability, good performance with heavy content such as multimedia and personalized experience with assertive recomendations.

\section{System Description}
Business problem, actor catalog, use case definitions, and
system boundary diagram



\subsection{Business Problem}
Netflix solves the problem of providing a large amount of audiovisual content to millions of users 
through the internet while giving each user a convenient and personalized experience.
\\Users want to watch movies and series whenever and wherever they want without having to download the entire content beforehand. 
At the same time, Netflix needs to help users find something they want to watch among a very large catalog.
\\The platform provide another services that not will be consider in this project like: Option for subtitles, interactive experiences or content preview.
 

\subsection{Actors}
The application has 4 users represented by a person, each of this person has a particular function and use the application for differents purposes. 
Also there are some system actor that provide services for the application. 

\begin{table}[H]
    \centering
    \renewcommand{\arraystretch}{1.3}
    
    \begin{tabular}{|
        >{\centering\arraybackslash}p{2.3cm}|
        p{4.3cm}|
        >{\centering\arraybackslash}p{2.3cm}|
        p{5.0cm}|
    }
        \hline
        \textbf{Actor} & 
        \textbf{Role} & 
        \textbf{Frequency of interaction} & 
        \textbf{Expectation} \\
        \hline

        Customer 
        & A person who has a monthly membership plan with the platform.
        & High
        & Be able to watch series and movies without interruptions and receive relevant recommendations.
        \\
        \hline

        Operational user
        & A member of the operations team responsible for managing content within the system.
        & Medium
        & Be able to manage all the content without problems.
        \\
        \hline

        Data analyst
        & A member of the Data \& Insights area who can access data and statistics.
        & Medium
        & Have access to useful data with clear statistics.
        \\
        \hline

        Platform administrator
        & Is responsible for managing users and the business model.
        & Low
        & Be able to manage the users of the system and membership configurations.
        \\
        \hline

        Media Convert service
        & It allows videos to be sent in small segments.
        & High
        & The user does not need to download the entire video; it is stored in fragments in the device buffer.
        \\
        \hline

        CDN service
        & Servers that help reduce latency by delivering content from the nearest server.
        & High
        & Ensure rapid response to high volumes of requests.
        \\
        \hline

        Storage
        & Storage for multimedia content and metadata, useful for performing queries.
        & High
        & Persist all data over time.
        \\
        \hline

    \end{tabular}

    \caption{\centering Actors, roles, interaction frequency, and expectations of the Netflix platform.}
    \label{tab:actors-netflix}
\end{table}

\subsection{Main Use Cases}
Subsection example 2

\subsection{System Boundaries}
Subsection example 2


\section{Functional Requirements}
This requriments specify what application does. There are 11 functional requirements divided in domain

\section{Non-Functional Requirements}
Numbered catalog of measurable non-functional
requirements per category

\section{Workload Model}
Structured table with all workload parameters and justification
for each value


\begin{thebibliography}{9}
\bibitem{rummery1994}
Rummery, G. A., \& Niranjan, M. (1994). On-line Q-learning using connectionist systems. \textit{Technical Report}, Cambridge University Engineering Department.


\end{thebibliography}

\end{document}